We thus showed that the specific geometry of the intersecting D-branes leads to different results when computing the value of the classical action, that is the leading contribution to the Yukawa couplings in string theory.
In particular in the Abelian case the value of the action is exactly the area formed by the intersecting D-branes in the $\R^2$ plane, that is the string worldsheet is completely contained in the polygon on the plane.
The difference between the \SO{4} case and \SU{2} is more subtle as in the latter there are complex coordinates in $\R^4$ for which the classical string solution is holomorphic in the upper half plane.
In the generic case presented so far this is in general no longer true.
The reason can probably be traced back to supersymmetry, even though we only dealt with the bosonic string.
In fact when considering \SU{2} rotated D-branes part of the spacetime supersymmetry is preserved, while this is not the case for \SO{4} rotations.

In the general case there does not seem to be any possible way of computing the action~\eqref{eq:action_with_imaginary_part} in term of the global data.
Most probably the value of the action is larger than in the holomorphic case since the string is no longer confined to a plane.
Given the nature of the rotation its worldsheet has to bend in order to be attached to the D-brane as pictorially shown in~\Cref{fig:brane3d} in the case of a $3$-dimensional space.
The general case we considered then differs from the known factorised case by an additional contribution in the on-shell action which can be intuitively understood as a small ``bump'' of the string worldsheet in proximity of the boundary.

In a technical and direct way we also showed the computation of amplitudes involving an arbitrary number of Abelian spin and matter fields.
The approach we introduced does not generally rely on \cft techniques and can be seen as an alternative to bosonization and old methods based on the Reggeon vertex.
Starting from this work the future direction may involve the generalisation to non Abelian spin fields and the application to twist fields.
In this sense this approach might be the only way to compute the amplitudes involving these complicated scenarios.
This analytical approach may also shed some light on the non existence of a technique similar to bosonisation for twist fields.


% vim: ft=tex
